\wordchapter{3}{Assessment of Government Intervention}

\wordsection{3.1}{Economic Rationale}

I assess each instrument by the constraint it addresses, the incentives it creates, and the distribution of its costs and benefits. Correlated losses motivate public risk sharing, while information and delivery costs motivate investment in insurance infrastructure (Miranda \& Glauber, 1997; Mahul \& Stutley, 2010). Premium assistance can serve affordability and distributional objectives (Hazell et al., 2017). Each instrument can affect production by changing prices, access, or risk bearing. Chapter 5 examines this possible consequence of insurance availability, complementing the assessment of program design.

\wordsection{3.2}{Premium Subsidies}

Let \(P\) denote the pure premium per unit of coverage and \(s\) the combined national and local subsidy rate. Policyholders then pay \((1 - s)P\) per unit. Given coverage demand \(D( \cdot )\), purchases without and with the subsidy are

\wordeq{0pt}{q_{0} = D(P),\quad\quad q_{1} = D\left( (1 - s)P \right)}{1}

In Figure 4, supply is perfectly elastic at \(P\). The subsidy increases coverage from \(q_{0}\) to \(q_{1}\) and costs the government \(sPq_{1}\). Of this amount, \(sPq_{0}\) supports purchases that would have occurred without the subsidy.

\begin{wordfigure}{4}{Premium Support and the Demand for Insurance Coverage}
\IfFileExists{"insurance graph.tex"}{%
  {\centering\adjustbox{max width=6in}{\input{"insurance graph"}}\par}%
}{%
  \wordimage{6in}{figures/fig4_premium_support.png}%
}
\wordnotes{\emph{Notes}: P is the pure premium per unit of coverage and s the combined subsidy rate. Supply is perfectly elastic at P, and the loading is financed separately. D is observed demand; D* is demand when the value of protection is fully perceived. The shaded rectangle is public pure-premium support, sPq₁, a transfer rather than a welfare loss. Arbitrary units.}
\end{wordfigure}

The figure holds the risk class and premium fixed and treats the loading as separately financed. The shaded rectangle measures a public transfer. If willingness to pay reflects the benefits of protection and P equals marginal resource cost, coverage above q₀ generates an allocative loss. If demand understates those benefits, D* shows why additional coverage might be warranted. The interpretation depends on these assumptions. Risk preferences, liquidity constraints, and information can affect demand differently, and changes in the insured risk pool can change premiums. The figure illustrates these mechanisms; it does not estimate welfare gains in Korea or identify an optimal subsidy rate.

Under the 2026 national schedule, support covers 33-65 percent of the pure premium for specified product-coverage combinations and 50 percent for other combinations. National support is capped at KRW 50 million per policyholder per product (MAFRA, 2026a, p. 2). At a given subsidy rate, the transfer rises with the eligible premium until the ceiling binds. Additional premiums then receive no further national subsidy. The ceiling limits transfer size but does not target household income, and local assistance remains a separate consideration.

\wordsection{3.3}{Administrative Cost Subsidies}

Administrative support finances the full loading under the 2026 rules (MAFRA, 2026a). The loading includes delegated selling and other specified operating costs. Loss-investigation expenses instead form part of the pure premium, together with the risk premium. Comparisons of administrative spending must keep these categories separate to avoid counting investigation costs twice. The relevant question is whether public financing delivers coverage and timely, accurate settlement at a reasonable resource cost.

Delivery costs can be assessed at two stages. Selling and policy-servicing expenditure can be measured per active contract, and claim-assessment expenditure per completed claim. Settlement times and assessment corrections provide corresponding measures of service. Operating-cost ratios can supplement these measures, provided that their premium denominator is clearly defined. Risk prices, claim frequency, product complexity, and geographical dispersion can change the ratios without changing efficiency. Comparisons should therefore account for these differences across products and regions.

Mahul and Stutley (2010, pp. 84--87) discuss the importance of low-cost delivery channels. Korea's 2026 guidelines also require payment within seven days after the indemnity has been determined and the policyholder submits a claim (MAFRA, 2026a, p. 5). This deadline applies after loss assessment; it does not measure the full time from loss notification to payment. Evaluating service therefore requires both compliance with existing rules and evidence on the cost, duration, and quality of the full process.

\wordsection{3.4}{State Reinsurance}

State reinsurance increases the capacity to cover correlated losses but creates a contingent public liability. Under a stop-loss contract with attachment point \(a\), total covered claims \(L\) are divided into retained and ceded claims as follows:

\wordeq{12pt}{L^{ret} = \min(L,a),\quad\quad L^{ced} = \max(L - a,0)}{2}

Under a proportional contract, the state receives a fixed share of premiums and pays the same share of claims. Korea used a stop-loss contract with a 180 percent attachment point in 2005--2013. Profit-and-loss sharing, which combines proportional and non-proportional reinsurance, began in 2014, extended to all crops in 2017, and replaced the previous arrangement fully in 2019 (MAFRA, 2025, pp. 267--268, 273--274). Figure 5 compares two benchmark contracts. Panel B uses a 50 percent state share for illustration; it does not reproduce the full Korean treaty.

In the uncapped stop-loss benchmark, private claims remain constant above the attachment point. Actual incentives also depend on coverage limits, pricing, and renewal terms. Proportional sharing retains private exposure at the margin but transfers part of ordinary losses to the state. The appropriate allocation depends on loss correlation and relative risk-bearing capacity (Jia et al., 2025). Let \(R\) denote the reinsurance premium in the same units as \(L\) and \(a\). Expected public cost under stop-loss equals expected ceded claims less the premium received:

\wordeq{12pt}{E\lbrack C\rbrack = E\left\lbrack \max(L - a,0) \right\rbrack - R}{3}

From 2005 to 2024, reported crop reinsurance payments totaled KRW 923.8 billion and premium receipts totaled KRW 470.0 billion (MAFRA, 2025, p. 269). The difference of KRW 453.8 billion is a cumulative nominal net payment, rather than a measure of actuarial adequacy or the fund's complete financial balance. Assessing adequacy requires comparing premiums with expected claims and the cost of bearing risk under the relevant treaties.

\begin{wordfigure}{5}{Allocation of Claims under Stop-Loss and Proportional Reinsurance}
\wordimage{6.2in}{figures/fig5_reinsurance_allocation.png}
\wordnotes{\emph{Notes}: Both axes express claims as percentages of premium. Panel A uses the 180 percent attachment point reported for 2005-2013 (MAFRA, 2025, p.~273). Private claims combine those of the operator and private reinsurers. In both panels, the private and state shares are stacked and sum to total claims. Panel B is an illustrative proportional contract without additional layers, premium transfers, expenses, or commissions.}
\end{wordfigure}

A prospective evaluation must distinguish expected net cost from liquidity needs and exposure to severe or consecutive loss years. Assessing each requires information on treaty terms, premiums, reserves, and committed financing. Catastrophe exposure depends on how losses co-move as well as on their frequency and severity (Mahul \& Stutley, 2010; Kang \& Chung, 2017). Historical payments should be assessed under the contracts that generated them, rather than under current sharing rates.
